\section*{अध्याय \ref{ch:b2:elements-in-the-earth} --- पृथ्वी में तत्त्व}

\begin{solution}{exo:b2:elements-in-the-earth:1}
पोटैशियम \ce{K2O} के द्रव्यमान का $2\times39.10/(2\times39.10+16.00) = 78.20/94.20 = 0.8301$ भाग है,
अतः $\num{5130}\times0.8301 = \qty{4259}{ppm}$ K। टाइटेनियम \ce{TiO2} का
$47.87/79.87 = 0.5993$ भाग है: $\num{27310}\times0.5993 = \qty{16370}{ppm}$ Ti।
\end{solution}

\begin{solution}{exo:b2:elements-in-the-earth:2}
ऐलुमिनियम शिलारागी और दुर्गलनीय है: सिलिकेट पृथ्वी ने उसे पूरा रखा, लगभग
कॉन्ड्राइटी अनुपात में (1.16)। सोडियम शिलारागी है पर वाष्पशील (\qty{958}{K}): वह
अभिवृद्धि से पहले या उसके दौरान खो गया, और 0.22 बचा। निकल मैग्नीशियम के साथ संघनित होता है
(\qty{1353}{K}), अतः वाष्पशीलता 0.077 की व्याख्या नहीं कर सकती: वह लौहरागी है, और उसका अधिकांश
क्रोड में चला गया। इरिडियम दुर्गलनीय (\qty{1603}{K}) और प्रबल लौहरागी है:
0.0029, लगभग सारा क्रोड में।
\end{solution}

\begin{solution}{exo:b2:elements-in-the-earth:3}
$\mathrm{La}_N = 31/0.232 = 133.6$, $\mathrm{Yb}_N = 2.0/0.163 = 12.27$, और
$(\mathrm{La}/\mathrm{Yb})_N = 10.9$। भूपर्पटी सभी दुर्लभ मृदाओं में समृद्ध है, और
भारी की तुलना में हल्की में ग्यारह गुना अधिक: वह गलितों से बनी है, जो
असंगत हल्की दुर्लभ मृदाओं को सबसे सहजता से ग्रहण करते हैं।
\end{solution}

\begin{solution}{exo:b2:elements-in-the-earth:4}
प्रसामान्यीकृत मान: $\mathrm{Sm}_N = 4.7/0.145 = 32.4$, $\mathrm{Eu}_N = 1.0/0.0546 = 18.3$,
$\mathrm{Gd}_N = 4.0/0.198 = 20.2$; $\mathrm{Eu}^* = \sqrt{32.4\times20.2} = 25.6$ और
$\mathrm{Eu}/\mathrm{Eu}^* = 18.3/25.6 = 0.72$, एक ऋणात्मक विसंगति। प्लेजियोक्लेज़, जो
\ce{Eu^2+} को ग्रहण करता है, उन गलितों से क्रिस्टलित हुआ जिन्होंने ऊपरी भूपर्पटी बनाई, और
नीचे गहराई में रह गया।
\end{solution}

\begin{solution}{exo:b2:elements-in-the-earth:5}
$\bar D_{\mathrm{La}} = 0.65\times0.058+0.18\times0.001+0.17\times0.001 = 0.0377 +
0.0002 + 0.0002 = 0.038$; $\bar D_{\mathrm{Lu}} = 0.3445+1.044+0.0306 = 1.42$। ल्यूटेशियम
संगत है ($\bar D > 1$), लैंथेनम प्रबल रूप से असंगत। गार्नेट चट्टान के ल्यूटेशियम का
$1.044/1.42 = 74$ प्रतिशत धारण करता है, यद्यपि वह चट्टान का केवल 18 प्रतिशत
है।
\end{solution}

\begin{solution}{exo:b2:elements-in-the-earth:6}
सोडियम के लिए, सिलिकेट पृथ्वी का Na/Mg $2670/\num{228000} = 0.01171$ है और CI
मान $5010/\num{95870} = 0.05226$: अनुपात 0.224 है। निकल के लिए, $1960/\num{228000}
= 0.008596$ बनाम $\num{10640}/\num{95870} = 0.1110$: 0.0775। सोडियम अपनी
वाष्पशीलता (\qty{958}{K}) के कारण न्यून है; निकल \qty{1300}{K} से ऊपर संघनित होता है और
सोडियम से तीन गुना अधिक न्यून है क्योंकि वह लौहरागी है: क्रोड ने उसका अधिकांश भाग
ले लिया।
\end{solution}

\begin{solution}{exo:b2:elements-in-the-earth:7}
La: $5.21\times0.058 = \qty{0.302}{ppm}$; Lu: $0.53\times0.53 = \qty{0.281}{ppm}$।
अनुपातों के अनुपात में प्रसामान्यीकरण-मान कट जाते हैं:
$(\mathrm{La}/\mathrm{Lu})_{N,\mathrm{cpx}}/(\mathrm{La}/\mathrm{Lu})_{N,\mathrm{melt}} =
D_{\mathrm{La}}/D_{\mathrm{Lu}} = 0.058/0.53 = 0.11$: क्रिस्टल का प्रतिरूप भारी दुर्लभ मृदाओं की ओर
गलित के प्रतिरूप से नौ गुना अधिक तीव्र है।
\end{solution}

\begin{solution}{exo:b2:elements-in-the-earth:8}
भूपर्पटी: $(12/0.658)/(31/0.648) = 18.24/47.84 = 0.38$। हवाई का बेसाल्ट:
$(18.10/0.658)/(15.20/0.648) = 27.51/23.46 = 1.17$। MORB: $7.96/8.04 = 0.99$। केवल
भूपर्पटी, नायोबियम की प्रबल कमी के साथ, क्षेपण-क्षेत्र का चिह्न धारण करती है।
\end{solution}

\begin{solution}{exo:b2:elements-in-the-earth:9}
\omterm{def:b2:elements-in-the-earth:bse}{आद्य प्रावार}: $240/0.0203 = \num{11800}$। MORB: K $= 0.160\times0.830\times10^4 =
\qty{1328}{ppm}$, K/U $= \num{11200}$। भूपर्पटी: K $= 2.80\times0.830\times10^4 =
\qty{23240}{ppm}$, K/U $= \num{8600}$। पोटैशियम प्रावार से भूपर्पटी तक सौ गुना बढ़ता है,
फिर भी K/U एक-तिहाई से कम बदलता है: दोनों तत्त्व लगभग समान रूप से असंगत हैं,
और गलन उन्हें मुश्किल से अलग करता है। यूरेनियम दुर्गलनीय है, इसलिए उसकी प्रचुरता
कॉन्ड्राइटी मॉडल से मिल जाती है; तब लावाओं पर मापा गया K/U ग्रह का पोटैशियम देता है, जिसे,
वाष्पशील होने के कारण, \omterm{def:b2:planet-formation-meteorites:chondrule}{कॉन्ड्राइटों} से नहीं लिया जा सकता।
\end{solution}

\begin{solution}{exo:b2:elements-in-the-earth:10}
साम्यावस्था पर हर तत्त्व के लिए $C_{\mathrm{rock}} = \bar D\,C_{\mathrm{melt}}$ (और
गलित का छोटा अंश चट्टान के संघटन को लगभग अपरिवर्तित छोड़ता है), अतः
$(C_1/C_2)_{\mathrm{melt}} = (C_{1}/\bar D_1)/(C_{2}/\bar D_2)$, चट्टान के $C_1,
C_2$ के लिए, और अनुपात $\bar D_2/\bar D_1$ है। La और Lu के लिए: $1.419/0.03805 = 37$। गार्नेट-युक्त
प्रावार का गलित अपने स्रोत की तुलना में Lu के सापेक्ष La में 37 गुना अधिक समृद्ध है: गार्नेट ने
भारी दुर्लभ मृदाओं को रोक रखा, और गलित का प्रतिरूप तीव्र ढलान वाला है।
\end{solution}

\begin{solution}{exo:b2:elements-in-the-earth:11}
हर तत्त्व के लिए $C_{\mathrm{mix}} = fC_U + (1-f)C_M$: $f$ में रैखिक। अनुपात
$\bigl(f\,\mathrm{Th}_U + (1-f)\mathrm{Th}_M\bigr)/\bigl(f\,\mathrm{Nb}_U +
(1-f)\mathrm{Nb}_M\bigr)$ दो रैखिक फलनों का भागफल है, $f$ में एक अतिपरवलय।
Th/Nb $= 0.2$ तब होता है जब $f(10.5-2.4) + (1-f)(0.404-1.048) = 0$, अर्थात् $8.1f = 0.644(1-f)$,
$f = 0.644/8.744 = 0.074$: लगभग 7 प्रतिशत भूपर्पटी Th/Nb को 0.077 से 0.2 तक उठा देती है।
$f = 0.1$ के लिए: Th $= 1.05 + 0.364 = 1.414$, Nb $= 1.2 + 4.716 = 5.916$, Th/Nb $= 0.24$।
\end{solution}

\begin{solution}{exo:b2:elements-in-the-earth:12}
किसी ऐसे घटक को हटाने से जिसमें दोनों में से कोई तत्त्व न हो, दोनों द्रव्यमान, और इसलिए उनका अनुपात,
अपरिवर्तित रहते हैं। यदि सारा ऐलुमिनियम रखा गया, तो $(\mathrm{Al}/\mathrm{Mg}) = 1.16$ गुना CI
का अर्थ है कि सिलिकेट पृथ्वी में अपने कॉन्ड्राइटी मैग्नीशियम का $1/1.16 = 0.86$ भाग है: 14 प्रतिशत
गायब है। ऐलुमिनियम के सापेक्ष सिलिकॉन CI मान का $0.83/1.16 = 0.72$ गुना है:
कॉन्ड्राइटी सिलिकॉन का 28 प्रतिशत गायब होता।
\end{solution}

\begin{solution}{pb:b2:elements-in-the-earth:1}
\textbf{1.} Fe/Mg $= \num{182800}/\num{95870} = 1.907$; Ni/Mg $= 0.1110$; Si/Mg $=
1.111$; Al/Mg $= 0.0887$।

\textbf{2.} परमाणु-अनुपात $(\num{182800}/55.85)/(\num{106500}/28.09) = 3273/3791 = 0.863$ लोहे के परमाणु प्रति
सिलिकॉन परमाणु, प्रकाशमंडल के 0.81 के 7 प्रतिशत के भीतर: सार्थक, क्योंकि चट्टान बनाने वाले तत्त्वों के लिए
\omterm{def:b2:elements-in-the-earth:ci}{CI कॉन्ड्राइट} सूर्य से मेल खाते हैं।

\textbf{3.} किसी ऐसे घटक की हानि जिसमें अनुपात का कोई तत्त्व न हो, अनुपात को
अपरिवर्तित छोड़ती है (\cref{exo:b2:elements-in-the-earth:12}): केवल निरपेक्ष सांद्रताएँ बदलती हैं।

\textbf{4.} मैग्नीशियम शिलारागी है और क्रोड में प्रवेश नहीं करता; वह एक मुख्य तत्त्व है,
जो \omterm{def:b2:planet-formation-meteorites:chondrule}{कॉन्ड्राइटों} और प्रावार दोनों में भली-भाँति निर्धारित है। लोहा वही तत्त्व है जो क्रोड और प्रावार
के बीच बँट रहा है: वह स्वयं अपना संदर्भ नहीं हो सकता।

\textbf{5.} जो तत्त्व साथ-साथ संघनित होते हैं वे नीहारिका में वाष्पशीलता से पृथक नहीं हुए,
इसलिए समग्र पृथ्वी को उनके अनुपात बनाए रखने चाहिए। Al/Mg हटता है: ऐलुमिनियम
दुर्गलनीय है (\qty{1653}{K}), और सिलिकेट पृथ्वी का Al/Mg CI मान का 1.16 गुना है।

\textbf{6.} Fe $= 8.05\times55.85/71.85 = 6.26$ प्रतिशत; Ni $= 0.25\times58.69/74.69 =
0.196$ प्रतिशत।

\textbf{7.} Fe/Mg $= 6.26/22.8 = 0.2746$; Ni/Mg $= 0.196/22.8 = 0.00860$।

\textbf{8.} लोहा: $1 - 0.2746/1.907 = 0.856$, 85.6 प्रतिशत गायब; निकल:
$1 - 0.00860/0.1110 = 0.923$, 92.3 प्रतिशत।

\textbf{9.} इरिडियम निकल से अधिक प्रबल लौहरागी है, और निकल लोहे से:
कोई तत्त्व धातु को जितना अधिक पसंद करता है, धातु उसका उतना ही अधिक भाग क्रोड में ले गई।

\textbf{10.} सिलिकेटों में, ओलिवीन और पाइरॉक्सीन में \ce{Fe^2+} के रूप में (8 प्रतिशत FeO):
FeO में ऑक्सीकृत लोहा \omterm{def:b2:elements-in-the-earth:goldschmidt}{शिलारागी तत्त्व} की तरह ऑक्सीजन से आबंधित होता है और
प्रावार में रह गया; केवल धात्विक लोहा डूबा।

\textbf{11.} $(1-x)C_{\mathrm{Fe}} + xc_{\mathrm{Fe}} = r_{\mathrm{Fe}}(1-x)C_{\mathrm{Mg}}$;
$(1-x)c_{\mathrm{Fe}}$ से भाग देने पर प्रतिज्ञप्ति मिलती है।

\textbf{12.} $r_{\mathrm{Fe}}C_{\mathrm{Mg}} - C_{\mathrm{Fe}} = 1.907\times22.8 - 6.26 =
37.22$; $y = 37.22/85.5 = 0.4353$, $x = y/(1+y) = 0.303$।

\textbf{13.} $0.1110\times22.8 - 0.196 = 2.335$; $y = 2.335/5.2 = 0.4490$, $x = 0.310$।

\textbf{14.} $y = (85.5\times37.22 + 5.2\times2.335)/(85.5^2 + 5.2^2) =
3194.4/7337.3 = 0.4354$, $x = 0.303$।

\textbf{15.} हर समीकरण का भार उसके गुणांक के वर्ग के बराबर होता है: लोहे के लिए $85.5^2/5.2^2 = 270$
गुना अधिक, जो कहीं अधिक प्रचुर तत्त्व है।

\textbf{16.} $37.22/2.335 = 15.9$, मॉडल के $85.5/5.2 = 16.4$ के सामने: 3 प्रतिशत
का अंतर।

\textbf{17.} $0.303\times\num{5.972e24} = \qty{1.81e24}{kg}$।

\textbf{18.} $(32.5-30.3)/32.5 = 7$ प्रतिशत कम।

\textbf{19.} $y = 0.325/0.675 = 0.4815$, क्रोड में Fe $= 37.22/0.4815 = 77.3$ प्रतिशत,
और Ni $= 2.335/0.4815 = 4.85$ प्रतिशत: 18 प्रतिशत हल्के तत्त्व, घनत्व द्वारा अनुमत 5 से
10 प्रतिशत का लगभग दोगुना। पृथ्वी का Fe/Mg CI मान से कुछ ऊँचा होना चाहिए,
या उसका सिलिकेट मैग्नीशियम मॉडल के मान से ऊँचा।

\textbf{20.} सिलिकॉन, सल्फ़र, कार्बन, ऑक्सीजन और हाइड्रोजन। सिलिकेट पृथ्वी में सिलिकॉन कम
है: उसका Mg से प्रसामान्यीकृत अनुपात, 0.83, दुर्गलनीय स्तर से नीचे है।

\textbf{21.} $1.111\times22.8 - 21.0 = 4.331$; $4.331/0.4354 = 9.9$ प्रतिशत Si
क्रोड में।

\textbf{22.} 9.9 बनाम 6 प्रतिशत: क्रोड में मॉडल से अधिक सिलिकॉन हो सकता है, या
सिलिकॉन का कुछ भाग कभी अभिवृद्ध ही नहीं हुआ। सिलिकॉन \qty{1310}{K} पर संघनित होता है, जो
मैग्नीशियम और लोहे (1336 और \qty{1334}{K}) से थोड़ा नीचे है: कुछ गैस में रह गया हो सकता है।

\textbf{23.} क्रोड: पृथ्वी के द्रव्यमान का $0.325\times0.5 = \qty{0.1625}{ppm}$; सिलिकेट
पृथ्वी: $0.675\times0.001 = \qty{0.000675}{ppm}$। क्रोड में सोने का
$0.1625/0.1632 = 99.6$ प्रतिशत है।

\textbf{24.} केवल रसायन पृथ्वी के द्रव्यमान का
\textbf{\qty[text-series-to-math=true]{30.3}{\percent}} उसके क्रोड में रखता है,
जो भूकंपविज्ञान के 32.5 प्रतिशत के 7 प्रतिशत के भीतर है।
\end{solution}
